\chapter{Keep the checksum state between reads}
\MARSChapterHeading{Keep the checksum state between reads}
\label{revision-lab:crc}

Use \texttt{crc32.asm}, \texttt{crc32-stream.asm} and the independent
\texttt{crc32-chunk.asm} in \texttt{revision/labs/files}. The main chapter
derives the reflected bit rule, table update and raw-state composition law.
This lab tests whether those distinctions survive actual instruction execution
and file reads. It uses ordinary Run and Debug; this unit's native guided lesson
remains unfinished.

Work with disposable source copies for faults. Each observation has a stated
boundary: a raw update return, an intermediate streaming state or a complete
file return. An exit syscall subsequently reuses \texttt{\$v0}; inspect the
intended boundary or the caller's stored result rather than treating the final
exit selector as a checksum.

\section{Predict the first byte before checking the whole message}

The default \texttt{crc32.asm} uses bit updates over nine ASCII bytes
\texttt{123456789}. The raw initial word is all ones. Before running, distinguish
selected count nine from its larger backing arena and from a NUL-terminated
string. A newline would be an additional input byte, not harmless formatting.

In a copy, change \texttt{input\_count} from nine to one. The selected byte is
ASCII \texttt{1}, hexadecimal \texttt{31}. XOR it into \texttt{ffffffff} to
obtain \texttt{ffffffce}. Predict the first two bit steps, writing both the old
low bit and shifted state before applying feedback. Your second result should
come from \texttt{3ffffff3 xor edb88320}, not from shifting the feedback mask.

Break at \texttt{bit\_next\_bit}. Step through the low-bit capture, logical
shift and conditional feedback. Compare your prediction with the main chapter's
eight-row trace. At \texttt{bit\_no\_feedback}, the decision for this step is
already complete, even if the branch did not skip an XOR. The old bit remains
in \texttt{\$t2}. After eight steps, expect raw state \texttt{7c231048}.
A count of one means one byte, not one bit step.

Restore count nine and break at \texttt{kernel\_done}. Expect raw
\texttt{340bc6d9}; \texttt{raw\_result} is still its initial $-1$ before the
caller's store. At \texttt{results\_ready}, its word holds that raw bit pattern.
Complementing the low 32 bits gives reported checksum \texttt{cbf43926}.
Explain why printing or storing the raw value under a label that promises the
reported checksum would be a contract mistake even if every bit step were right.

For diagnosis, change only the bit loop's \texttt{dsrl} to the 32-bit arithmetic
\texttt{sra} in a fresh copy. Find the first differing state. Bit 31 and the
upper half of the register matter. A full-width \texttt{dsra} is a different
experiment: it agrees with the logical shift while bit 63 remains zero. State
which width you tested rather than declaring every arithmetic shift equivalent
or every arithmetic shift incorrect.

\section{Continue the returned state through another method}

Open the unchanged \texttt{crc32-stream.asm}. Its demonstration caller is
specifically arranged for the nine-byte fixture: four bytes through the bit
routine, then five through the table routine. The reusable update functions
retain their ordinary pointer/count interfaces. Changing this demonstration's
input length also requires changing its two slice lengths; it is not a general
file reader.

Break at \texttt{stream\_first}. Predict the role of the returned value before
reading it: it describes the first four bytes and has not been finalized.
Expect \texttt{641c1f5c} in \texttt{\$v0}. The next instruction retains it in
\texttt{\$s0} across table construction. The caller then supplies the next
source address, count five and that retained raw state to the table update.

\begin{figure}[htbp]
\centering
\begin{tikzpicture}[alt={Two actual light-theme Debug observations of the complete stream example. After four bytes v0 is hexadecimal 00000000641c1f5c. After continuing through the other five bytes it is 00000000340bc6d9. Both full-register displays retain zero upper halves. Native labels identify prefix and continued raw state; neither value is finalized.}]
\node[anchor=north west,inner sep=0] (a) at (0,0) {\MARSRevisionCrop{revision/captures/crc/crc-prefix-raw.png}{0.47\linewidth}{0.24375}{0.33111111}{0.59375}{0.63777778}};
\node[anchor=north west,inner sep=0] (b) at ([xshift=4mm]a.north east) {\MARSRevisionCrop{revision/captures/crc/crc-complete-raw.png}{0.47\linewidth}{0.24375}{0.33111111}{0.59375}{0.63777778}};
\node[anchor=south west,inner sep=0,font=\small\bfseries] at ([yshift=2mm]a.north west) {After four bytes};
\node[anchor=south west,inner sep=0,font=\small\bfseries] at ([yshift=2mm]b.north west) {After nine bytes};
\begin{scope}[shift={(a.south west)},x={($(a.south east)-(a.south west)$)},y={($(a.north west)-(a.south west)$)}]
\draw[black,line width=.7pt] (.03,.12) rectangle (.97,.90);
\end{scope}
\begin{scope}[shift={(b.south west)},x={($(b.south east)-(b.south west)$)},y={($(b.north west)-(b.south west)$)}]
\draw[black,line width=.7pt] (.03,.12) rectangle (.97,.90);
\end{scope}
\end{tikzpicture}
\caption{Actual raw return values at the two boundaries of the supplied stream example, observed in separate Debug executions. Continuing the returned prefix state produces the same nine-byte result as one whole update. Native LaTeX supplies crops, labels and boxes.}
\label{fig:revision-crc-native-states}
\end{figure}


At \texttt{stream\_done}, expect \texttt{340bc6d9}. This is the same state
as the single nine-byte update, despite the buffer split and method change.
The native observations use hexadecimal display in full 64-bit registers;
leading zeroes express the required zero extension.

In separate copies, try each of two faults: reset the second seed to
\texttt{ffffffff}, or complement the first result before passing it as that
seed. Predict why neither operation preserves the original message. Find the
first changed interface value, rather than comparing only the final checksum.
An empty update should preserve an arbitrary incoming state; it does not mean
that every empty operation returns the finalized empty-message checksum zero.

\section{Prepare a table, then account for its cost}

In an unchanged \texttt{crc32.asm}, choose exercise one and break at
\texttt{kernel\_done}. Expect return value 256, meaning entries initialized.
Inspect table words at byte offsets zero, four, 512 and 1020. They hold
\texttt{00000000}, \texttt{77073096}, \texttt{edb88320} and
\texttt{2d02ef8d}. The words after the selected 1024-byte table retain their
sentinel values. Distinguish the entry index 128 from its byte offset 512.

A fresh execution restores initial data, including the unprepared table. Merely
changing exercise to two and starting over does not preserve the table from a
previous execution. Use the complete stream caller, which prepares it and uses
it in one session, or the behavior checks, which explicitly supply prepared
entries. A lookup into allocated sentinel words is not an alternative CRC table.

Repeat the chapter's measurement with the complete source:
\begin{verbatim}
java -jar /path/to/cli.jar grade \
  revision/labs/checks/crc32-work.kts \
  revision/labs/files/crc32.asm
\end{verbatim}
The \texttt{INIT} row reports table construction; \texttt{METRIC} rows report
input name, byte count, method, completed instructions and finalized host-oracle
checksum. Method zero is bit updating and one is table updating. The checker
verifies behavior as well as recording work. Count construction once when your
comparison requires it.

For the 256-byte case, compare 14668 bit instructions with $14348+3331$ for a
new table and one update. Then repeat the reasoning for 4096 bytes. Explain
which setup boundary reverses the apparent choice. Do not translate these
ratios into elapsed times. If you change a loop, retain the original input,
measure the new source and run correctness checks again.

\section{Checksum a real file, including the previous image}

From the public book repository, run the actual host-file helper with new output
directories. Replace the Java and CLI paths:
\begin{verbatim}
python3 tools/run_revision_crc.py \
  --java /path/to/java --cli /path/to/cli.jar \
  --output /tmp/crc-bit
python3 tools/run_revision_crc.py \
  --java /path/to/java --cli /path/to/cli.jar \
  --method table --output /tmp/crc-table
\end{verbatim}
Both default runs write the nine-byte fixture, execute the ordinary file syscall
path and require finalized checksum \texttt{cbf43926} with status zero. They
obtain the assembled result addresses in a separate nonexecuting grading query,
then inspect a memory dump after the real file run. No grading file fixture is
installed for that run. Each directory retains source, input, dump and receipt;
the helper refuses an existing output directory.

Use \texttt{--bytes 8193 --capacity 13 --method table} for a generated binary
file spanning many reads. The generator includes high-byte values and zeroes.
Use \texttt{--bytes 0} for an empty file. Predict a successful empty file's
checksum and explain why its status still needs to be checked.

Finally use \texttt{--input /path/to/output.ppm} with your Chapter 18 output.
This checks every byte of that file, including its header. For the unchanged
23-byte red/green/blue/white fixture, expect \texttt{e481fe30}. A semantically
identical image with a different accepted header can have a different whole-file
CRC because the selected byte sequence changed. Retain the artifact and receipt
so that your comparison names the exact file.

\section{Distinguish EOF from failure after useful work}

The native failure observation below uses a separate authoring fixture with
nine input bytes, capacity seven, maximum read three and an injected error
after five accepted bytes. Its actual descriptor is 19. Two positive reads
return three and two; the next read returns $-1$, followed by a close.
The state after that prefix exists, but the function must report failure to
read the complete file.

\begin{figure}[htbp]
\centering
\begin{tikzpicture}[alt={Actual light-theme Debug return after a controlled read failure. Register v0 is zero and v1 is two in decimal display. The separate virtual-file receipt records descriptor nineteen, positive reads three and two, a failed read, and close. The image establishes the return values; it does not itself display descriptor cleanup.}]
\node[anchor=north west,inner sep=0] (a) {\MARSRevisionCrop{revision/captures/crc/crc-file-error.png}{0.23\linewidth}{0.2453125}{0.31333333}{0.69921875}{0.63777778}};
\begin{scope}[shift={(a.south west)},x={($(a.south east)-(a.south west)$)},y={($(a.north west)-(a.south west)$)}]
\draw[black,line width=.7pt] (.04,.02) rectangle (.89,.98);
\end{scope}
\end{tikzpicture}
\caption{The actual file return reports zero checksum and read-failure status two. It does not publish a successful checksum for the five-byte prefix. Controlled I/O and close evidence are retained separately; native LaTeX supplies the crop and outline.}
\label{fig:revision-crc-native-error}
\end{figure}


At \texttt{file\_done}, expect checksum zero and status two, before the caller
stores either word. The native operation receipt verifies the read sequence and
closed descriptor; the register image shows only the return values. This is
controlled virtual I/O, separate from the actual host-file runs above.

Run the public file cases with descriptor 19:
\begin{verbatim}
java -jar /path/to/cli.jar grade \
  revision/labs/checks/crc32-files.kts \
  revision/labs/files/crc32.asm
\end{verbatim}
They test missing files, EOF, short reads and injected failures while checking
workspace and saved state. Correct buffering may request fewer than capacity
bytes; it must use the actual returned descriptor and process only the actual
positive return count. A short positive read is not EOF.

Draw a timeline for success and another for a later read failure. Mark the one
initialization, each raw-state update, the close, and finalization where it is
allowed. Explain how resetting state per read, processing scratch capacity after
a short read, or reporting a prefix as success would change the result or hide
a failure. These are separate faults even when they all produce an unexpected
checksum on one convenient input.

\section{Write the bounded chunk checker}

Complete \texttt{crc32\_check\_chunk} in a copy of
\texttt{revision/labs/files/crc32-chunk.asm}. The main chapter defines the
pointer/count/table interface and return statuses. Preserve its provided helpers
and demonstration caller; use either update strategy. Its selected buffer is a
single chunk, not a whole file, and its maximum selected count is 4096 bytes.

First prove the selected count is from twelve through 4096 before pointer access.
Then recover the unsigned big-endian length from four bytes and compare it with
selected count minus twelve. Reject disagreement before reading other fields or
the table. On a proved extent, checksum type plus payload, finalize once and
compare with the stored big-endian value. Return the computed checksum with
status zero for a match or one for a mismatch. Extent failure returns zero and
status two. Keep all caller bytes and full saved state unchanged.

The default twelve-byte fixture is an empty \texttt{IEND} chunk. Its covered
input is still four type bytes, and its correct checksum is \texttt{ae426082}.
Changing the last stored byte should produce status one while retaining that
computed value. Changing its encoded payload length to one without extending
the selected buffer should instead produce zero/status two before CRC work.
The original \texttt{break} stub intentionally fails these requirements.

\begin{verbatim}
java -jar /path/to/cli.jar grade \
  revision/labs/checks/crc32-chunk.kts \
  /path/to/your/crc32-chunk.asm
\end{verbatim}
The checks accept equivalent implementations. They distinguish empty payloads,
high bytes, offsets, mismatches, unsigned length errors and forbidden reads;
they do not certify every possible chunk or a complete PNG. Finish with a short
account of exactly which bytes your routine covers and where it first proves
the address of the stored CRC. That proof will be useful when a decoder begins
following a sequence of records rather than one supplied interval.
